% Propietario conservado: /Users/ruben/Documents/New project/output/REVISION_ARGUMENTAL_APERTURAS_CIERRES_20260915/III/spanish_source/manuscrito/sections/04c_velocidad_y_reloj.tex
% RESULTADO_RECUPERADO: funtor dimensional integral, eq:cZ-internos,
% eq:longitud-ruta; composición reunida en Artículo V REV02, sección 69.
% Incorporación autónoma a III: mismo transporte, mismas hojas y misma carta.
\subsection{Velocidad constitutiva y realización cinemática del reloj}
\label{iii:iii:sec:enlace-velocidad-constitutiva}

La respuesta constitutiva y el lector cinemático realizan una misma
sección de velocidad. Su identificación conserva el orden de construcción:
las hojas de APP, su orientación TRIT y el transporte angular publicado
por el TPK anteceden a las lecturas de velocidad y a la elección de unidades.
En la representación de dos hojas de la
sección~\ref{iii:iii:sec:hojas}, con $x>|y|$ y proyectores de rango uno
$P_\pm$, el transporte es
\[
 T=e^{-xI-y\mathcal R}=q_+P_++q_-P_-,
 \qquad q_\pm=e^{-x\mp y},\qquad 0<q_\pm<1.
\]
Sus bloques temporales de longitudes $90$ y $120$ producen la respuesta
$\mathcal V$ de \eqref{iii:iii:eq:vacuum-op}. Las coordenadas $x,y$ y los
canales son salidas del desarrollo angular previo; esta composición
aplica sus lectores a ese transporte ya construido.

El lector logarítmico simétrico del transporte temporal es
\begin{equation}
 \mathcal E=
 \log\!\bigl[(1-q_+^{90})(1-q_-^{90})\bigr]
 -\log\!\bigl[(1-q_+^{120})(1-q_-^{120})\bigr].
 \label{iii:iii:eq:velocidad-lector-simetrico}
\end{equation}
En la recta dimensional de velocidad, con sección positiva $u_C$, el
lector cinemático publica $c_{\mathrm{int}}=\exp(-\mathcal E)u_C$.
La respuesta constitutiva publica
$c_{\mathrm{vac}}=\widehat c u_C$ mediante \eqref{iii:iii:eq:four}.

\begin{proposition}[Identidad constitutiva y cinemática]
\label{iii:iii:prop:identidad-velocidad}
En la misma carta dimensional se verifica
\begin{equation}
 c_{\mathrm{int}}=c_{\mathrm{vac}}=\widehat c\,u_C,
 \qquad \widehat c=(r_+r_-)^{-1}=\exp(-\mathcal E).
 \label{iii:iii:eq:identidad-velocidad}
\end{equation}
La identificación conserva el intercambio de hojas.
\end{proposition}
\begin{proof}
Todos los factores de \eqref{iii:iii:eq:velocidad-lector-simetrico} son
positivos. La ley del logaritmo y \eqref{iii:iii:eq:rchannels} dan
\[
 \mathcal E=\log(r_+r_-)=\log\det\mathcal V,
 \qquad e^{-\mathcal E}=(\det\mathcal V)^{-1}.
\]
El determinante corresponde a la realización de dos hojas de rango uno.
Las dos definiciones publican, por tanto, la misma sección. El intercambio
de hojas permuta $r_+$ y $r_-$ y conserva su producto. El determinante
$\det T=e^{-2x}$ corresponde al transporte elemental, mientras que
$\det\mathcal V=r_+r_-$ corresponde a su respuesta temporal;
\eqref{iii:iii:eq:identidad-velocidad} utiliza el segundo.
\end{proof}

\subsubsection{Longitud de trayectoria y duración elemental}

El marco dimensional positivo $(u_S,u_C,u_T,u_Q)$ induce la base de
longitud $u_L=u_Cu_T$. Para una trayectoria enriquecida $\gamma$ de
$n=|\gamma|$ transiciones y canal constante, los lectores aditivos son
\begin{equation}
 T_{\mathrm{int}}(\gamma)=nu_T,
 \qquad L_{\mathrm{int}}(\gamma)=n\widehat c\,u_L.
 \label{iii:iii:eq:velocidad-longitud-reloj}
\end{equation}
La concatenación suma el número de transiciones y sus realizaciones.
Las rutas enriquecidas conservan, además, su orientación y su memoria.

\begin{corollary}[Normalización de la longitud elemental]
Para $n>0$, $L_{\mathrm{int}}(\gamma)/T_{\mathrm{int}}(\gamma)
=c_{\mathrm{vac}}$. En la normalización $t_0=u_T$,
\begin{equation}
 \ell_0=c_{\mathrm{int}}t_0=\widehat c\,u_L.
 \label{iii:iii:eq:longitud-elemental-constitutiva}
\end{equation}
En una carta temporal $t_0=\tau_0u_T$, con $\tau_0>0$, se obtiene
$\ell_0=\tau_0\widehat c\,u_L$.
\end{corollary}
\begin{proof}
Dividir las dos expresiones de \eqref{iii:iii:eq:velocidad-longitud-reloj}
cancela $n$ y utiliza $u_L/u_T=u_C$. Multiplicar la sección de velocidad
por $t_0$ prueba ambas fórmulas elementales.
\end{proof}

Cuando el canal depende de la arista, el lector conserva la forma aditiva
$L_{\mathrm{int}}(\gamma)=\sum_{e\in\gamma}\widehat c(e)u_L$.
Su cociente por $nu_T$ es la media de las velocidades de esas aristas.
La expresión de canal constante es la especialización homogénea de
esta lectura.

\subsubsection{Cambio de carta y normalización adaptada}

\begin{proposition}[Covariancia de la identificación]
Sean $u_C'=d u_C$ y $u_T'=\eta u_T$, con $d,\eta>0$.
Las coordenadas de las mismas secciones satisfacen
\begin{equation}
 \widehat c'=d^{-1}\widehat c,
 \qquad \tau_0'=\eta^{-1}\tau_0,
 \qquad u_L'=d\eta u_L.
 \label{iii:iii:eq:velocidad-cambio-carta}
\end{equation}
Si otra publicación escribe $c'=\widehat c' u_C'$ con $u_C'=d u_C$,
la igualdad de secciones equivale exactamente a $d\widehat c'=\widehat c$.
\end{proposition}
\begin{proof}
La igualdad de cada sección en las dos bases da
$\widehat c'u_C'=\widehat c u_C$ y $\tau_0'u_T'=\tau_0u_T$.
En particular,
\[
 \tau_0'\widehat c'u_L'
 =\frac{\tau_0}{\eta}\frac{\widehat c}{d}\,d\eta u_L
 =\tau_0\widehat c u_L=\ell_0.
\]
La última equivalencia procede de la misma sustitución para la velocidad.
\end{proof}

La elección adaptada $d=\widehat c$ da $u_C'=c_{\mathrm{vac}}$ y
$\widehat c'=1$: expresa la sección construida en una base adaptada.
El coeficiente de la carta interna sigue siendo $(r_+r_-)^{-1}$, y esa
es la normalización utilizada por la inversión cúbica. Manteniendo las
bases de acción y carga, \eqref{iii:iii:eq:dimensions} produce
\[
 \widehat\mu'=d\widehat\mu,\qquad
 \widehat\varepsilon'=d\widehat\varepsilon,\qquad
 \widehat\mu'\widehat\varepsilon'=(\widehat c')^{-2}.
\]
En la carta adaptada se tiene $\widehat\mu'=r_-/r_+$ y
$\widehat\varepsilon'=r_+/r_-$, por sustitución de
$\widehat\mu=r_-^2$ y $\widehat\varepsilon=r_+^2$.

\subsubsection{Conservación del lector bajo amplificación}

Para el refinamiento $\mathcal V_m=\mathcal V\otimes I_{9^m}$, la
traza normalizada $\tau_m=\operatorname{Tr}/(2\cdot9^m)$ satisface
\begin{equation}
 \exp[-2\tau_m(\log\mathcal V_m)]
 =\exp[-\log r_+-\log r_-]=\widehat c.
 \label{iii:iii:eq:velocidad-traza-normalizada}
\end{equation}
En efecto, cada autovalor aparece $9^m$ veces; esa multiplicidad se
cancela con el denominador de la traza normalizada. El determinante
ordinario ampliado es $(r_+r_-)^{9^m}$. La normalización de dos hojas,
equivalentemente expresada por \eqref{iii:iii:eq:velocidad-traza-normalizada},
conserva el lector constitutivo durante la amplificación. Esta
realización conserva el transporte enriquecido que la origina.
